\documentclass{situationist-preprint}

\title{Analyticity and Borderline Cases}
\msid{LNG-04-b}
\msfamily{LNG-04}
\msdiscipline{Philosophy of Language}
\mstype{Companion manuscript}
\mskeywords{analyticity; vagueness; borderline cases; Quine; meaning change; revisability}

\begin{document}
\maketitle

\begin{abstract}
LNG-04-a argues that the analytic/synthetic distinction is vague, not empty: it has clear positive cases, clear negative cases and borderline cases. This manuscript examines the borderline cases. We identify their source in sentences whose rejection would be neither clearly a revision of belief nor clearly a change of meaning, classify them, and show that their existence cannot be turned into an argument against the clear cases. Quine's argument from revisability, read charitably, is a sorites on \enquote{analytic}; like any sorites it fails at the inductive step, and it no more empties the distinction than the sorites empties \enquote{heap}.
\end{abstract}

\section{The source of borderline cases}

LNG-04-a defines a sentence as analytic if its acceptance is licensed by the rules that constitute the meanings of its words (Definition~2.2), and holds that constitution comes in degrees (Principle~P1): some rules governing a word are central to its meaning, others peripheral, and there is no sharp line between them. Proposition~3.3 there infers that there are borderline cases. We make the inference explicit.

\begin{definition}\label{d:revmean}
A speaker who comes to reject a sentence she previously accepted \emph{revises a belief} if her words continue to mean what they meant; she \emph{changes a meaning} if at least one word comes to mean something different.
\end{definition}

\begin{proposition}\label{p:source}
A sentence is a borderline case of analyticity if and only if its rejection by a competent community would be neither clearly a revision of belief nor clearly a change of meaning.
\end{proposition}

\begin{proof}[Argument]
If acceptance of a sentence is clearly licensed by meaning-constituting rules, rejecting it clearly changes what some word means, since a rule constitutive of the word is abandoned; so its rejection is clearly a change of meaning. If acceptance is clearly not so licensed, rejecting it leaves the constitutive rules in place, and is clearly a revision of belief. Conversely, where neither verdict is clear, the rule licensing acceptance is neither clearly central nor clearly peripheral, and by Principle~P1 of LNG-04-a the sentence is neither clearly analytic nor clearly synthetic.
\end{proof}

\section{Three kinds of borderline case}

\emph{Law-cluster terms.} Putnam observed that terms such as \enquote{energy} or \enquote{momentum} are governed by a cluster of laws, none of which is definitional; when \enquote{kinetic energy equals one half of mass times velocity squared} was given up in relativistic mechanics, it is unclear whether the word changed meaning or a belief was corrected \citep{Putnam1962}. The answer is unclear because the abandoned law was central to the use of the term but not uniquely so.

\emph{Natural kind terms.} \enquote{Whales are fish} was accepted, and its rejection was a discovery. \enquote{Cats are animals} seems analytic, yet Putnam imagined that cats might turn out to be robots controlled from Mars \citep{Putnam1962}. On the account of LNG-08-a, the reference-fixing profile of a kind term includes a sortal under which its samples are taken (Definition~2.4 there), and rules tied to the sortal are central to the term's use; a sentence that states the sortal (\enquote{gold is a metal}) is close to analytic, while a sentence that states a superficial classification (\enquote{whales are fish}) is not. Sentences linking a kind term to a higher kind lie in between, according to how firmly the sortal of the original grounding is fixed.

\emph{Open-textured everyday terms.} Whether \enquote{a person is a human being} is analytic depends on whether the rule linking \enquote{person} to \enquote{human} is constitutive of the word or is a generalization ordinary speakers accept. Disputes about persons who are not human have the character of disputes in which both sides are competent, which is the mark of a borderline case.

In each kind, the sentences at issue are those that the clear cases of LNG-04-a (\enquote{all bachelors are unmarried}, \enquote{if $p$ and $q$ then $p$}; Proposition~3.1) are not. A stipulative definition introduces a rule that is constitutive by fiat; the rule for conjunction is the one whose abandonment would leave a connective with different truth-conditions, as LNG-01-a shows (Theorem~3.3). No comparable argument is available for \enquote{whales are mammals} or for any physical law.

\section{The revisability sorites}

Quine's second argument is that any statement can be held true come what may, and that no statement is immune from revision \citep{Quine1951}. Read as a challenge to the distinction, it runs as follows. Some statements once thought analytic, such as the law of excluded middle or the claim that the shortest path between two points is a straight line, have been revised or proposed for revision. Between these and the clearest cases there is a series of statements each of which is as revisable as its neighbor. So the clearest cases are revisable too, and nothing is analytic.

This is a sorites. Its inductive premise, that if a statement in the series is revisable without change of meaning so is its neighbor, is the tolerance principle for \enquote{revisable without change of meaning}. That predicate is vague by Proposition~\ref{p:source}, and tolerance principles for vague predicates are false: whatever the correct account of vagueness, the inductive premise of a sorites fails somewhere in the series. On the epistemic account defended in LNG-12, it fails at a sharp point we cannot locate \citep{Williamson1994}; on other accounts it fails without a sharp point. On either, the conclusion that every statement is revisable without change of meaning does not follow, and the clear cases remain clear. What the argument establishes is what LNG-04-a concedes: that there are statements whose status is not settled.

\begin{objection}
A vague distinction is useless for the purposes for which analyticity was introduced, such as explaining a priori knowledge.
\end{objection}

\begin{reply}
A distinction with clear cases can do explanatory work in the clear cases. That knowledge of \enquote{bachelors are unmarried} needs no empirical support is explained by its analyticity, whatever we say about \enquote{momentum is mass times velocity}. Most useful distinctions, between living and nonliving or between child and adult, are vague in the same way.
\end{reply}

\begin{objection}
Williamson argues that for any sentence, competent speakers can reject it while understanding it, through idiosyncratic theoretical views; so understanding never requires assent \citep{Williamson2007}.
\end{objection}

\begin{reply}
Williamson's target is the view that understanding a sentence requires being disposed to assent to it. The definition used here is different: a sentence is analytic if its acceptance is licensed by meaning-constituting rules, whether or not every competent speaker follows the license. A logician who rejects an instance of a classical law on theoretical grounds is either using the connective with different rules, and so with a different meaning, or making an error about what the rules license. Which description applies to a given case may be unclear, and where it is, the case is borderline.
\end{reply}

\section{Conclusion}

Borderline cases of analyticity arise where the line between revising a belief and changing a meaning is unclear. They are real and numerous, but they cannot be used to spread unclarity to the clear cases, since the argument that would do so is a sorites. The distinction is vague, and there are analytic truths.

\begin{thebibliography}{9}\small
\bibitem[Putnam(1962)]{Putnam1962} Putnam, H. (1962). The analytic and the synthetic. In H.~Feigl and G.~Maxwell (eds.), \emph{Minnesota Studies in the Philosophy of Science}, Vol.~3. University of Minnesota Press.
\bibitem[Quine(1951)]{Quine1951} Quine, W.~V.~O. (1951). Two dogmas of empiricism. \emph{Philosophical Review}, 60, 20--43.
\bibitem[Williamson(1994)]{Williamson1994} Williamson, T. (1994). \emph{Vagueness}. Routledge.
\bibitem[Williamson(2007)]{Williamson2007} Williamson, T. (2007). \emph{The Philosophy of Philosophy}. Blackwell.
\end{thebibliography}

\end{document}
